% GENERATED FILE. Edit canonical lessons, then run npm run generate:books.
% canonical-source-hash: fnv1a64:a51937b673bba9f9
% canonical-lessons: UR-C217-darhi, UR-C217-munchh, UR-C217-tandurusti, UR-C217-dant-dard, UR-C217-khujli

\chapter{Beard, Moustache, Good health, Toothache, Itch}
\label{ch:ur-beard-moustache-good-health-toothache-itch}
\hlchaptermodality{217}

\begin{chapteropening}
\textbf{By the end of this chapter:} \emph{I can say a beard, a moustache, good health, a toothache and an itch.}

Complete the last lesson of chapter 217: I can say a beard, a moustache, good health, a toothache and an itch.
\end{chapteropening}

\section[\texorpdfstring{\ur{داڑھی} (dāṛhī)}{داڑھی (dāṛhī)}]{\ur{داڑھی} (dāṛhī) — a beard}

\label{lesson:UR-C217-darhi}



\begin{quote}
Before the new one: say the Urdu for a dream, then the Urdu for weakness.
\end{quote}

\subsection*{You'll want to know: \ur{داڑھی}}
\textbf{\ur{داڑھی}} — \emph{dāṛhī} — \textquotedblleft{}a beard\textquotedblright{}.

\begin{cousinweb}
\textbf{\ur{داڑھی}} (\emph{dāṛhī}) is a tadbhava, a word inherited and reshaped by sound change. It comes from \emph{*dāṃṣṭra}, \textquotedblleft{}fang, tusk, beard\textquotedblright{}, a form CDIAL reconstructs from its descendants (CDIAL 6250). It is built on Sanskrit \emph{daṃṣṭrā}, \textquotedblleft{}tusk, fang\textquotedblright{}; CDIAL finds a change of sense from \textquotedblleft{}fang\textquotedblright{} through \textquotedblleft{}tusk\textquotedblright{} to \textquotedblleft{}beard\textquotedblright{} possible.
\end{cousinweb}

\begin{grammarlens}[title={What you've built}]
One more word for this chapter.
\end{grammarlens}

\subsection*{Your turn}
\begin{itemize}
\raggedright
  \item \emph{Say it:} \emph{dāṛhī}
  \item \emph{Say it:} \emph{dāṛhī}, once more
  \item \emph{Say it:} say \emph{dāṛhī}
  \item \emph{Recall:} say the Urdu for a dream, then the Urdu for weakness, then say \emph{dāṛhī} again
\end{itemize}

\begin{tcolorbox}[breakable,colback=teal!4,colframe=teal!35!black,title={Before you move on}]
What does \ur{داڑھی} mean? (\textquotedblleft{}A beard\textquotedblright{}.) Say it once more.
\end{tcolorbox}
\section[\texorpdfstring{\ur{مونچھ} (mūṅchh)}{مونچھ (mūṅchh)}]{\ur{مونچھ} (mūṅchh) — a moustache}

\label{lesson:UR-C217-munchh}



\begin{quote}
Before the new one: say the Urdu for weakness, then the Urdu for a beard.
\end{quote}

\subsection*{You'll want to know: \ur{مونچھ}}
\textbf{\ur{مونچھ}} — \emph{mūṅchh} — \textquotedblleft{}a moustache\textquotedblright{}.

\begin{cousinweb}
\textbf{\ur{مونچھ}} (\emph{mūṅchh}) is a tadbhava, a word inherited and reshaped by sound change. It comes from \emph{*mucchā}, \textquotedblleft{}moustache\textquotedblright{}, a form CDIAL reconstructs from its descendants (CDIAL 10180).
\end{cousinweb}

\begin{grammarlens}[title={What you've built}]
One more word for this chapter.
\end{grammarlens}

\subsection*{Your turn}
\begin{itemize}
\raggedright
  \item \emph{Say it:} \emph{mūṅchh}
  \item \emph{Say it:} \emph{mūṅchh}, once more
  \item \emph{Say it:} say \emph{mūṅchh}
  \item \emph{Recall:} say the Urdu for weakness, then the Urdu for a beard, then say \emph{mūṅchh} again
\end{itemize}

\begin{tcolorbox}[breakable,colback=teal!4,colframe=teal!35!black,title={Before you move on}]
What does \ur{مونچھ} mean? (\textquotedblleft{}A moustache\textquotedblright{}.) Say it once more.
\end{tcolorbox}
\section[\texorpdfstring{\ur{تندرستی} (tandurustī)}{تندرستی (tandurustī)}]{\ur{تندرستی} (tandurustī) — good health}

\label{lesson:UR-C217-tandurusti}



\begin{quote}
Before the new one: say the Urdu for a beard, then the Urdu for a moustache.
\end{quote}

\subsection*{You'll want to know: \ur{تندرستی}}
\textbf{\ur{تندرستی}} — \emph{tandurustī} — \textquotedblleft{}good health\textquotedblright{}.

\begin{cousinweb}
\textbf{\ur{تندرستی}} (\emph{tandurustī}) is \emph{tandurust}, \textquotedblleft{}healthy\textquotedblright{}, usually traced to Persian \emph{tandurust}, \textquotedblleft{}sound of body\textquotedblright{}, with the ending \emph{-ī} that makes a noun.
\end{cousinweb}

\begin{grammarlens}[title={What you've built}]
One more word for this chapter.
\end{grammarlens}

\subsection*{Your turn}
\begin{itemize}
\raggedright
  \item \emph{Say it:} \emph{tandurustī}
  \item \emph{Say it:} \emph{tandurustī}, once more
  \item \emph{Say it:} say \emph{tandurustī}
  \item \emph{Recall:} say the Urdu for a beard, then the Urdu for a moustache, then say \emph{tandurustī} again
\end{itemize}

\begin{tcolorbox}[breakable,colback=teal!4,colframe=teal!35!black,title={Before you move on}]
What does \ur{تندرستی} mean? (\textquotedblleft{}Good health\textquotedblright{}.) Say it once more.
\end{tcolorbox}
\section[\texorpdfstring{\ur{دانت} \ur{درد} (dāṅt dard)}{دانت درد (dāṅt dard)}]{\ur{دانت} \ur{درد} (dāṅt dard) — a toothache}

\label{lesson:UR-C217-dant-dard}



\begin{quote}
Before the new one: say the Urdu for a moustache, then the Urdu for good health.
\end{quote}

\subsection*{You'll want to know: \ur{دانت} \ur{درد}}
\textbf{\ur{دانت} \ur{درد}} — \emph{dāṅt dard} — \textquotedblleft{}a toothache\textquotedblright{}.

\begin{cousinweb}
\textbf{\ur{دانت} \ur{درد}} (\emph{dāṅt dard}) joins \emph{dāṅt}, \textquotedblleft{}tooth\textquotedblright{}, from Sanskrit \emph{danta}, \textquotedblleft{}tooth\textquotedblright{} (CDIAL 6152), and \emph{dard}, \textquotedblleft{}pain\textquotedblright{}, usually traced to Persian \emph{dard}.
\end{cousinweb}

\begin{grammarlens}[title={What you've built}]
One more word for this chapter.
\end{grammarlens}

\subsection*{Your turn}
\begin{itemize}
\raggedright
  \item \emph{Say it:} \emph{dāṅt dard}
  \item \emph{Say it:} \emph{dāṅt dard}, once more
  \item \emph{Say it:} say \emph{dāṅt dard}
  \item \emph{Recall:} say the Urdu for a moustache, then the Urdu for good health, then say \emph{dāṅt dard} again
\end{itemize}

\begin{tcolorbox}[breakable,colback=teal!4,colframe=teal!35!black,title={Before you move on}]
What does \ur{دانت} \ur{درد} mean? (\textquotedblleft{}A toothache\textquotedblright{}.) Say it once more.
\end{tcolorbox}
\section[\texorpdfstring{\ur{کھجلی} (khujlī)}{کھجلی (khujlī)}]{\ur{کھجلی} (khujlī) — an itch}

\label{lesson:UR-C217-khujli}



\begin{quote}
Before the new one: say the Urdu for good health, then the Urdu for a toothache.
\end{quote}

\subsection*{You'll want to know: \ur{کھجلی}}
\textbf{\ur{کھجلی}} — \emph{khujlī} — \textquotedblleft{}an itch\textquotedblright{}.

\begin{cousinweb}
\textbf{\ur{کھجلی}} (\emph{khujlī}) is a tadbhava, a word inherited from Sanskrit and reshaped by sound change. It comes from \emph{kharju}, \textquotedblleft{}itching, scab\textquotedblright{} (CDIAL 3827). CDIAL takes the word from Dravidian. The Sanskrit word is known only from the old word-lists.
\end{cousinweb}

\begin{grammarlens}[title={What you've built}]
One more word for this chapter.
\end{grammarlens}

\subsection*{Your turn}
\begin{itemize}
\raggedright
  \item \emph{Say it:} \emph{khujlī}
  \item \emph{Say it:} \emph{khujlī}, once more
  \item \emph{Say it:} say \emph{khujlī}
  \item \emph{Recall:} say the Urdu for good health, then the Urdu for a toothache, then say \emph{khujlī} again
\end{itemize}

\begin{tcolorbox}[breakable,colback=teal!4,colframe=teal!35!black,title={Before you move on}]
What does \ur{کھجلی} mean? (\textquotedblleft{}An itch\textquotedblright{}.) Say it once more.
\end{tcolorbox}
